\subsection{第二存在定理}

定理 3. 设 A 为整闭整环，A' 为包含 A 且在 A 上整的环。假设 A 中在 A' 里是零因子的元素只有 0。设 p、q 为 A 的两个素理想，满足 q ⊃ p，并设 q' 为 A' 的位于 q 之上的一个素理想。那么存在 A' 的位于 p 之上的素理想 p'，使 q' ⊃ p'。

先设 A' 是整环。设 K、K' 分别为 A 与 A' 的分式域；设 K'' 为 K' 的代数闭包，A'' 为 A 在 K'' 中的整闭包；于是 \(A \subset A' \subset A''\)。设 \(p''\) 为 A'' 的位于 p 之上的素理想（第 1 节定理 1），\(q''\) 为 A'' 的位于 q 之上且满足 \(p'' \subset q''\) 的素理想（第 1 节定理 1 的推论 2），最后设 \(q_{1}''\) 为 A'' 的位于 \(q'\) 之上的素理想（第 1 节定理 1）。由第 3 节命题 6 (i)，存在 K'' 的一个 K-自同构 \(\sigma\) 使 \(\sigma \cdot q'' = q_{1}''\)。于是 \(\sigma \cdot p''\) 是 A'' 的位于 p 之上且满足 \(\sigma \cdot p'' \subset q_{1}''\) 的素理想，从而 \(p' = A' \cap \sigma \cdot p''\) 是 A' 的位于 p 之上且包含于 \(A' \cap q_{1}'' = q'\) 的素理想。

现在转到一般情形。由于 A 是整环且 \(q'\) 是素理想，子集 \(A - \{0\}\) 与 \(A' - q'\)（皆为 \(A'\) 的子集）都是乘性的；于是它们的积 \(S = (A - \{0\})(A' - q')\) 是 \(A'\) 的一个不含 0 的乘性子集，因为 A 的非零元素都不是 \(A'\) 中的零因子。于是存在（第 II 章 §2 第 5 节命题 11 的推论 2）一个素理想 \(m'\)（属于 \(A'\)），它与 S 不相交，换言之，它满足 \(m' \subset q'\) 且 \(m' \cap A = 0\)。设 h 为典范同态 \(A' \to A'/m'\)。h 在 A 上的限制是单射，因而 \(h(A)\) 是整闭的。由于 \(m' \subset q'\)，\(h(q')\) 是 \(A'/m'\) 的位于 \(h(q)\) 之上的素理想；因为 \(A'/m'\) 是整环，证明的第一部分表明存在一个素理想 \(n'\)（属于 \(A'/m'\)），使 \(n' \cap h(A) = h(p)\) 且 \(h(q') \supset n'\)。理想 \(p' = h^{-1}(n')\) 是 \(A'\) 的素理想，且 \(q' \supset p'\)，因为 \(q'\) 包含 h 的核。由于 \(n' \supset h(p)\)，有 \(p' \supset p\)。最后，对 \(x \in p' \cap A\)，有 \(h(x) \in n' \cap h(A) = h(p)\)，又因 h 在 A 上的限制是单射，故 \(x \in p\)；因此 \(p' \cap A = p\)。

推论. 在定理 3 关于 A 与 A' 的假设下，设 p 为 A 的一个素理想。A' 的位于 p 之上的素理想就是 A' 中包含 pA' 的素理想所成之集 E 的极小元。

A' 的位于 p 之上的素理想依据第 1 节命题 1 的推论 1 是 E 中的极小元。反之，设 q' 为 E 的一个极小元。由于 q' ∩ A ⊃ p，定理 3 表明存在 A' 的位于 p 之上的素理想 p' 使 q' ⊃ p'。由于 p' 包含 pA'，对 q' 所作的假设蕴含 q' = p'，从而 q' 位于 p 之上。

* 设 V、V' 为两个仿射代数簇，f 为 V' 到 V 的一个态射，使得 \(f(V')\) 在 V 中稠密。设 A（相应地 A'）为 V（相应地 V'）上正则函数的环；给定了 f 之后，我们可以把 A 等同于 A' 的一个子环；设 A' 在 A 上整。第 1 节的定理 1 表明，对 V 的每个不可约子簇 W，存在 \(W'\)，即 \(V'\) 的一个不可约子簇，使 \(f(W')\) 是 W 的稠密子集；特别地，V 的每个点都是 \(V'\) 的某个不可约子簇的像，这表明 f 是满射。类似地，f 在每个不可约子簇 \(W'\)（\(V'\) 的）上的限制把 \(W'\) 映到 V 的一个不可约子簇上。第 1 节定理 1 的推论 2 表明：若 W 与 X 是 V 的两个不可约子簇，满足 \(W \supset X\)，且若 \(W'\) 是 \(V'\) 的一个满足 \(f(W') = W\) 的不可约子簇，则存在不可约子簇 \(X'\)（\(V'\) 的、包含于 \(W'\) 的），使 \(f(X') = X\)。

若 A 是整闭的，则称 V 为正规的。定理 3 表明：若 V 正规，若 W 与 X 是 V 的两个不可约子簇，满足 \(W \supset X\)，且若 \(X'\) 是 \(V'\) 的一个满足 \(f(\mathbf{X}') = \mathbf{X}\) 的不可约子簇，则存在子簇 \(W'\)（在 \(V'\) 中），它包含 \(X'\) 且满足 \(f(W') = W\)。最后，定理 3 的推论表明：若 V 正规且 W 是 V 的一个不可约子簇，则不可约子簇 \(W'\)（属于 \(V'\)、满足 \(f(W') = W\) 者）恰是 \(f(W)\) 的不可约分支。

